% 正文主结果的补充过程图；不改变各问正式费用口径。
\section{补充结果图}\label{sec:app_figures}
\renewcommand{\thesubsection}{\Alph{section}.\arabic{subsection}}

正文各问只保留主表、主图和结论数字。本附录补充调度过程、费用路径、净负荷时空结构和波动电价分解，用来核对这些结论是否与数据机制一致，不另给一套正式结果。

\Needspace{8\baselineskip}
\subsection{问题一：充放电是否跟着电价走}

问题一的指定时段表已给出各 $4$ 小时充放电量。图~\ref{fig:app_q1_cd} 把同一套线性规划方案画成与分时电价对照的时序，用来核对“低价充电、高价放电”是否在典型日上成立，而不是只看加总后的表。

\begin{figure}[!htbp]
  \centering
  \includegraphics[width=0.90\linewidth]{Q1_charge_discharge.png}
  \caption{典型日储能充放电与分时电价对照：低价时段充电、高峰放电}
  \label{fig:app_q1_cd}
\end{figure}

由图~\ref{fig:app_q1_cd} 可见，充电集中在电价低谷（约 $0{:}00$--$2{:}00$、正午及 $22{:}00$ 附近），放电集中在早晚高峰。$4{:}00$--$8{:}00$ 放电 $6\,608$ kWh、$16{:}00$--$20{:}00$ 放电 $5\,632$ kWh，与正文表~\ref{tab:q1_storage} 一致；正午低价窗口充电 $5\,507$ kWh，对应表中 $12{:}00$--$16{:}00$ 的充电。电价曲线与柱状方向相反，说明费用下降来自跨时段搬运，而不是改变典型日的负荷或光伏。全天弃电量为 $0$，日末储电量回到 $6000$ kWh，与正文约束一致。

\Needspace{8\baselineskip}
\subsection{问题二：回测区间费用从哪几天来}

正文问题二给出全年总费用 $13\,695\,403.61$ 元及策略对照。图~\ref{fig:app_q2_daily} 把该正式方案拆到 $2025$ 年 $2$ 月 $1$ 日至 $12$ 月 $31$ 日共 $334$ 天，用来看费用是均匀分摊，还是由少数高价日和紧急购电尖峰抬高。图~\ref{fig:app_q2_month} 再按月压缩，用来看季节而不是单日噪声。

\begin{figure}[!htbp]
  \centering
  \includegraphics[width=0.92\linewidth]{Q2_fig1_daily_cost.png}
  \caption{问题二正式方案逐日购电费：星期周期明显，紧急购电呈稀疏尖峰}
  \label{fig:app_q2_daily}
\end{figure}

由图~\ref{fig:app_q2_daily} 可见，日总费用呈稳定的星期起伏，工作日明显高于周末，与附件~2 负载的星期周期一致。计划购电费构成大多数日期的费用主体；紧急购电只在少数日期形成尖峰，说明日前锁定计划在多数日子可以执行，高价补购风险集中在供需预报偏差较大的几天。因此，问题二的分位修正和日末留电，主要是为压住这些稀疏尖峰，而不是改写每一天的计划购电。

\begin{figure}[!htbp]
  \centering
  \includegraphics[width=0.86\linewidth]{q2_monthly_cost_boxplot.pdf}
  \caption{问题二正式方案分月日费用：春末夏初较低，十二月最高}
  \label{fig:app_q2_month}
\end{figure}

由图~\ref{fig:app_q2_month} 可见，$2$--$5$ 月日费用中位随光伏增强下降，$6$--$7$ 月因负荷升高回升，$12$ 月中位和上须均最高。箱线上下跨度表明，同一月份内不同日期的费用可相差数万元，全年合计会把这种日间波动抹平。分月箱线用于核对正文全年费用是否被冬季高价日和夏季个别紧急购电日共同抬高，而不能单独看成“某月策略失效”。

\Needspace{8\baselineskip}
\subsection{净负荷何时富余、何时缺口}

正文预处理用日电量和持续曲线说明：全年 $80.7\%$ 时段净负荷为正，且正负两端都超过储能接口 $5000$ kW。图~\ref{fig:app_netload} 把同一净负荷按日期和时刻展开，用来看富余和缺口分别发生在哪些时段、哪些季节，从而解释为何必须同时设置弃光和紧急购电。

\begin{figure}[!htbp]
  \centering
  \includegraphics[width=0.92\linewidth]{Q2_eda_10_netload_heatmap.png}
  \caption{全年净负荷热图：正午富余集中在夏半年，早晚缺口全年都在}
  \label{fig:app_netload}
\end{figure}

由图~\ref{fig:app_netload} 可见，蓝色负净负荷主要出现在 $4$--$10$ 月的 $10{:}00$--$15{:}00$，即光伏超过负荷、需要充电或弃光的窗口；红色正净负荷在早晚高峰全年都在，冬季白天也多为正值。热图说明储能的充电窗口短而集中、放电窗口长而重复，接口功率约束会在两端同时起作用：夏半年正午可能充不满全部富余，早晚高峰也可能放不足全部缺口。这与持续曲线给出的极值结论一致，但补上了“发生在何时”。

\Needspace{8\baselineskip}
\subsection{问题三：夏季代表日能量如何平衡}

正文图~\ref{fig:q3_dispatch} 给出 $2025$ 年 $6$ 月 $21$ 日的负荷、光伏、购电功率和储电量。图~\ref{fig:app_q3_balance} 把同一日画成能量平衡堆叠，用来核对接力关系：负荷由哪些电量满足，弃光出现在什么时候，晚高峰靠什么削峰。费用与正文指定日结果一致，为 $22\,107.38$ 元，当日无紧急购电。

\begin{figure}[!htbp]
  \centering
  \includegraphics[width=0.92\linewidth]{fig_q3_balance_0621.pdf}
  \caption{夏季代表日能量平衡：正午充满后弃光，晚高峰由储能放电削峰}
  \label{fig:app_q3_balance}
\end{figure}

由图~\ref{fig:app_q3_balance} 可见，午间光伏峰值约 $8606$ kW，储电量触及上限后出现弃光，瞬时约 $4630$ kW，全天弃光 $11\,684$ kWh；晚高峰约 $19{:}00$ 由放电削峰约 $4005$ kW。堆叠满足负荷 $+$ 充电 $=$ 光伏利用 $+$ 调整购电 $+$ 放电。与正文折线图对照：折线看购电和储电量何时升降，本图看弃光和充放电在功率上如何接上。该日无紧急购电，说明夏季光伏窗口足以把电留到晚高峰，但不能消纳正午全部富余，因此弃光变量不能省。

\Needspace{8\baselineskip}
\subsection{问题四：费用升高是电价变了，还是预报错了}

正文问题四给出 Q4-2 全年 $14\,413\,479.02$ 元、Q4-3 的 LA 全年 $14\,001\,199.26$ 元。与固定分时电价下的问题二、问题三相比，费用都升高。图~\ref{fig:app_q4_decomp} 把升高拆成两段：实时电价环境差额，以及电价预测误差的代价，用来判断问题四变贵的主因。

\begin{figure}[!htbp]
  \centering
  \includegraphics[width=0.96\linewidth]{q4_error_decomposition.png}
  \caption{波动电价费用分解：环境差额约占 $4.6\%$，电价预测误差不足 $0.7\%$}
  \label{fig:app_q4_decomp}
\end{figure}

由图~\ref{fig:app_q4_decomp} 可见，若调度时使用真实电价（Oracle），Q4-2 和 Q4-3 的费用相对原固定电价情景分别升高约 $4.64\%$ 和 $4.76\%$，对应约 $63.6$ 万元和 $63.2$ 万元；在此之上，用 $0{:}00$ 电价预测做决策再按真实电价结算，只再增加约 $8.3$ 万元和 $9.3$ 万元，相对升幅约 $0.58\%$ 和 $0.67\%$。因此，问题四费用高于前两问，主要是附件 $4$ 的实时电价整体更贵，而不是电价预报把购电计划带偏。Q4-3 的实际费用 $1400.1$ 万元仍低于 Q4-2 的 $1441.3$ 万元，与正文“日内滚动仍能降低紧急购电”一致。
